%% kapa: exam=1 round=33 year=2022 period=1 form=A subject=real_estate subject_ko=부동산학원론 seq=12 number=92
%% source: Q-Net 감정평가사 기출문제 (cst00309) artlSeq=5209997 file=raw/1008/5209997/2022년도 제33회 감정평가사 1차시험 1교시(민법,경제학,부동산학).pdf
%% conversion: kordoc 4.24.51 -> tools/gen1.py (auto); review=proofread-sonnet; flags=none
\begin{question}{92}
정부의 주택시장 개입에 관한 설명으로 옳지 않은 것은?
\begin{choices}
  \item 주택은 긍정적인 외부효과를 창출하므로 생산과 소비를 장려해야 할 가치재(merit goods)이다.
  \item 저소득층에 대한 임대주택 공급은 소득의 직접분배효과가 있다.
  \item 주택구입능력을 제고하기 위한 정책은 소득계층에 따라 달라진다.
  \item 자가주택 보유를 촉진하는 정책은 중산층 형성과 사회안정에 기여한다.
  \item 주거안정은 노동생산성과 지역사회에 대한 주민참여를 제고하는 효과가 있다.
\end{choices}
\end{question}
\answer{2}
